\documentclass[11pt]{article}

\usepackage[a4paper,margin=1in]{geometry}
\usepackage{amsmath,amssymb,amsthm,mathtools}
\usepackage{microtype}

\newtheorem{theorem}{Theorem}
\newtheorem{proposition}{Proposition}
\newtheorem{lemma}{Lemma}
\newtheorem{corollary}{Corollary}
\newtheorem{example}{Example}
\newtheorem{remark}{Remark}

\DeclareMathOperator{\Tr}{Tr}
\DeclareMathOperator{\Cov}{Cov}
\DeclareMathOperator{\Ent}{Ent}
\DeclareMathOperator{\Id}{Id}

\newcommand{\R}{\mathbb R}
\newcommand{\cPtwo}{\mathcal P_2}
\newcommand{\cG}{\mathcal G}
\newcommand{\bSplus}{\mathbb S_+}
\newcommand{\W}{W}
\newcommand{\N}{\mathsf N}
\newcommand{\dd}{\,\mathrm d}

\title{Gelbrich's inequality and Gaussian-preserving Wasserstein flows}
\author{}
\date{}

\begin{document}

\maketitle

Gelbrich's inequality gives a sharp lower bound on the quadratic Wasserstein distance in terms of only means and covariance matrices. It was proved by Gelbrich in \cite{Gelbrich1990}. If
\[
m_\mu:=\int_{\R^d}x\,\dd\mu(x),
\qquad
\Sigma_\mu:=\int_{\R^d}(x-m_\mu)(x-m_\mu)^{\mathsf T}\,\dd\mu(x),
\]
then, for $\mu,\nu\in\mathcal P_2(\R^d)$,
\[
\W_2^2(\mu,\nu)
\ge
\|m_\mu-m_\nu\|^2
+
\Tr \Sigma_\mu+\Tr \Sigma_\nu
-
2\Tr\!\left((\Sigma_\mu^{1/2}\Sigma_\nu\Sigma_\mu^{1/2})^{1/2}\right).
\]
The right-hand side is precisely the squared $\W_2$ distance between the Gaussian laws with the same means and covariances as $\mu$ and $\nu$.

\section{Definitions}

Let $\mathcal P_2(\R^d)$ be the space of Borel probability measures on $\R^d$ with finite second moment. For $\mu,\nu\in\mathcal P_2(\R^d)$,
\[
\W_2^2(\mu,\nu)
:=
\inf_{\pi\in\Pi(\mu,\nu)}
\int_{\R^d\times\R^d}\|x-y\|^2\,\dd\pi(x,y),
\]
where $\Pi(\mu,\nu)$ denotes the set of couplings of $\mu$ and $\nu$.

For $A,B\in\mathbb S_+^d$, define
\[
\mathfrak B^2(A,B)
:=
\Tr A+\Tr B
-
2\Tr\!\left((A^{1/2}BA^{1/2})^{1/2}\right).
\]
This is the squared Bures distance between covariance matrices.

Let
\[
\mathcal G_d:=\{\N(m,A):m\in\R^d,\ A\in\mathbb S_+^d\}
\]
be the set of possibly degenerate Gaussian laws. We write
\[
R\mu:=\N(m_\mu,\Sigma_\mu)
\]
for the Gaussian law with the same mean and covariance as $\mu$.

\section{Gelbrich's inequality}

\begin{theorem}[Gelbrich's inequality]
Let $\mu,\nu\in\mathcal P_2(\R^d)$. Then
\[
\W_2^2(\mu,\nu)
\ge
\|m_\mu-m_\nu\|^2+\mathfrak B^2(\Sigma_\mu,\Sigma_\nu).
\]
Equivalently,
\[
\W_2^2(\mu,\nu)
\ge
\|m_\mu-m_\nu\|^2
+
\Tr \Sigma_\mu+\Tr \Sigma_\nu
-
2\Tr\!\left((\Sigma_\mu^{1/2}\Sigma_\nu\Sigma_\mu^{1/2})^{1/2}\right).
\]
\end{theorem}

\begin{proof}
We first prove the matrix estimate needed below.

\begin{lemma}
Let $X,Y$ be centered square-integrable $\R^d$-valued random variables with
\[
\Cov(X)=A,
\qquad
\Cov(Y)=B,
\]
where $A,B\in\mathbb S_+^d$. Then
\[
\Tr \mathbb E[XY^{\mathsf T}]
\le
\Tr\!\left((A^{1/2}BA^{1/2})^{1/2}\right).
\]
\end{lemma}

\begin{proof}
Set
\[
C:=\mathbb E[XY^{\mathsf T}].
\]
The block covariance matrix
\[
\begin{pmatrix}
A & C\\
C^{\mathsf T} & B
\end{pmatrix}
\]
is positive semidefinite.

Assume first that $A$ and $B$ are positive definite. Let
\[
K:=A^{-1/2}CB^{-1/2}.
\]
By congruence,
\[
\begin{pmatrix}
I & K\\
K^{\mathsf T} & I
\end{pmatrix}\succeq 0.
\]
Hence $\|K\|_{\mathrm{op}}\le 1$. By duality between the operator norm and the nuclear norm,
\[
\Tr C
=
\Tr(A^{1/2}KB^{1/2})
=
\Tr(KB^{1/2}A^{1/2})
\le
\|B^{1/2}A^{1/2}\|_*.
\]
Moreover,
\[
\|B^{1/2}A^{1/2}\|_*
=
\Tr\!\left((A^{1/2}BA^{1/2})^{1/2}\right).
\]
Thus the claim holds when $A,B$ are positive definite.

For general $A,B\in\mathbb S_+^d$, apply the positive-definite case to $A+\varepsilon I$ and $B+\varepsilon I$, then let $\varepsilon\downarrow 0$.
\end{proof}

Let $\pi\in\Pi(\mu,\nu)$, and let $(U,V)$ have law $\pi$. Define
\[
X:=U-m_\mu,
\qquad
Y:=V-m_\nu.
\]
Then
\[
\Cov(X)=\Sigma_\mu,
\qquad
\Cov(Y)=\Sigma_\nu.
\]
Also,
\[
\mathbb E\|U-V\|^2
=
\|m_\mu-m_\nu\|^2+\mathbb E\|X-Y\|^2.
\]
Expanding,
\[
\mathbb E\|X-Y\|^2
=
\Tr\Sigma_\mu+\Tr\Sigma_\nu
-
2\Tr\mathbb E[XY^{\mathsf T}].
\]
By the lemma,
\[
\Tr\mathbb E[XY^{\mathsf T}]
\le
\Tr\!\left((\Sigma_\mu^{1/2}\Sigma_\nu\Sigma_\mu^{1/2})^{1/2}\right).
\]
Therefore every coupling $\pi$ satisfies
\[
\int\|u-v\|^2\,\dd\pi(u,v)
\ge
\|m_\mu-m_\nu\|^2+\mathfrak B^2(\Sigma_\mu,\Sigma_\nu).
\]
Taking the infimum over $\pi\in\Pi(\mu,\nu)$ gives the result.
\end{proof}

\begin{corollary}
The map $R:\mathcal P_2(\R^d)\to\mathcal G_d$ is $1$-Lipschitz:
\[
\W_2(R\mu,R\nu)\le \W_2(\mu,\nu).
\]
\end{corollary}

\begin{proof}
For Gaussian laws,
\[
\W_2^2(\N(m,A),\N(n,B))
=
\|m-n\|^2+\mathfrak B^2(A,B).
\]
Applying Gelbrich's inequality to $\mu$ and $\nu$ gives
\[
\W_2^2(R\mu,R\nu)
=
\|m_\mu-m_\nu\|^2+\mathfrak B^2(\Sigma_\mu,\Sigma_\nu)
\le
\W_2^2(\mu,\nu).
\]
\end{proof}

\section{A Gaussian-preservation criterion}

The following criterion was communicated by Hugo Lavenant.

\begin{proposition}[Variational Gaussian-preservation criterion]
Let $F:\mathcal P_2(\R^d)\to(-\infty,+\infty]$ satisfy
\[
F(R\mu)\le F(\mu)
\qquad
\text{for every }\mu\in\mathcal P_2(\R^d).
\]
Fix $\tau>0$, and consider one step of the JKO scheme:
\[
\nu\in\operatorname*{argmin}_{\eta\in\mathcal P_2(\R^d)}
\left\{
F(\eta)+\frac{1}{2\tau}\W_2^2(\gamma,\eta)
\right\}.
\]
If $\gamma\in\mathcal G_d$, then $R\nu$ is also a minimizer. In particular, if the minimizer is unique, then $\nu=R\nu$ and $\nu$ is Gaussian.

Consequently, whenever the Wasserstein gradient flow of $F$ is unique and obtained as the limit of the JKO scheme, the flow preserves Gaussian initial data.
\end{proposition}

\begin{proof}
Since $\gamma$ is Gaussian,
\[
R\gamma=\gamma.
\]
For any competitor $\eta$,
\[
F(R\eta)\le F(\eta)
\]
by assumption, and
\[
\W_2(R\gamma,R\eta)\le \W_2(\gamma,\eta)
\]
by the previous corollary. Hence
\[
F(R\eta)+\frac{1}{2\tau}\W_2^2(\gamma,R\eta)
\le
F(\eta)+\frac{1}{2\tau}\W_2^2(\gamma,\eta).
\]
Applying this to a minimizer $\eta=\nu$ shows that $R\nu$ is also a minimizer. If the minimizer is unique, then $\nu=R\nu$, hence $\nu$ is Gaussian.
\end{proof}

\section{Examples satisfying the criterion}

\begin{example}[Functionals depending only on mean and covariance]
Let
\[
F(\mu)=\Phi(m_\mu,\Sigma_\mu)
\]
for some function
\[
\Phi:\R^d\times\mathbb S_+^d\to(-\infty,+\infty].
\]
Then
\[
F(R\mu)=F(\mu).
\]
Thus the JKO scheme, when started from a Gaussian law, admits a Gaussian minimizer at each step.

In particular, quadratic interaction energies are covered. If
\[
F(\mu)
=
\frac12\int_{\R^d\times\R^d}(x-y)^{\mathsf T}Q(x-y)\,\dd\mu(x)\dd\mu(y),
\]
with $Q=Q^{\mathsf T}$, then
\[
F(\mu)=\Tr(Q\Sigma_\mu).
\]
Hence $F$ depends only on $\Sigma_\mu$, so $F(R\mu)=F(\mu)$.
\end{example}

\begin{proof}
The first claim is immediate from
\[
m_{R\mu}=m_\mu,
\qquad
\Sigma_{R\mu}=\Sigma_\mu.
\]
For the quadratic interaction,
\[
\mathbb E[(X-Y)^{\mathsf T}Q(X-Y)]
=
\Tr\!\left(Q\,\mathbb E[(X-Y)(X-Y)^{\mathsf T}]\right),
\]
where $X,Y$ are independent with law $\mu$. Since
\[
\mathbb E[(X-Y)(X-Y)^{\mathsf T}]=2\Sigma_\mu,
\]
the prefactor $1/2$ gives
\[
F(\mu)=\Tr(Q\Sigma_\mu).
\]
\end{proof}

\begin{example}[Boltzmann entropy]
Let
\[
\Ent(\mu)
:=
\begin{cases}
\displaystyle \int_{\R^d}\rho\log\rho\,\dd x,
& \mu=\rho\,\dd x,\\[0.8em]
+\infty,
& \text{otherwise}.
\end{cases}
\]
Then
\[
\Ent(R\mu)\le \Ent(\mu).
\]
Hence the entropy flow preserves Gaussian initial data.
\end{example}

\begin{proof}
Assume first that $\mu=\rho\,\dd x$ has finite entropy and nondegenerate covariance $\Sigma_\mu$. Let
\[
g
\]
be the density of $R\mu=\N(m_\mu,\Sigma_\mu)$. Since relative entropy is nonnegative,
\[
0\le \int_{\R^d}\rho\log\frac{\rho}{g}\,\dd x
=
\int_{\R^d}\rho\log\rho\,\dd x
-
\int_{\R^d}\rho\log g\,\dd x.
\]
Now $\log g$ is a constant plus a quadratic polynomial:
\[
\log g(x)
=
c-\frac12(x-m_\mu)^{\mathsf T}\Sigma_\mu^{-1}(x-m_\mu).
\]
Because $\mu$ and $R\mu$ have the same mean and covariance,
\[
\int_{\R^d}\rho\log g\,\dd x
=
\int_{\R^d}g\log g\,\dd x
=
\Ent(R\mu).
\]
Therefore
\[
\Ent(\mu)\ge \Ent(R\mu).
\]
The singular case follows by the usual lower-semicontinuous extension, or by restricting to the affine span of the covariance.
\end{proof}

\begin{example}[Squared Wasserstein distance to a Gaussian target]
Let $\gamma\in\mathcal G_d$ and define
\[
F(\mu):=\frac12\W_2^2(\mu,\gamma).
\]
Then
\[
F(R\mu)\le F(\mu).
\]
\end{example}

\begin{proof}
Since $\gamma$ is Gaussian,
\[
R\gamma=\gamma.
\]
Therefore, by the $1$-Lipschitz property of $R$,
\[
\W_2(R\mu,\gamma)
=
\W_2(R\mu,R\gamma)
\le
\W_2(\mu,\gamma).
\]
Squaring gives $F(R\mu)\le F(\mu)$.
\end{proof}



\begin{thebibliography}{9}

\bibitem{Gelbrich1990}
M.~Gelbrich,
\newblock On a formula for the $L^2$-Wasserstein metric between measures on Euclidean and Hilbert spaces,
\newblock \emph{Mathematische Nachrichten} \textbf{147} (1990), 185--203.

\bibitem{LavenantPC}
H.~Lavenant,
\newblock Gaussian-preserving Wasserstein flows,
\newblock personal communication, October 2025.

\end{thebibliography}

\end{document}